\documentclass[letterpaper,12pt]{article}

\usepackage{titlesec}
\usepackage{fancyhdr}
\usepackage[most]{tcolorbox}
\usepackage{longtable}
\usepackage{makecell}
\usepackage{tikz}
\usepackage{geometry}
\usepackage{array}
\usepackage{multicol}
\usepackage[table]{xcolor}
\usepackage{polyglossia}
\setmainlanguage{english}
\setotherlanguage{arabic}
\setotherlanguage{hebrew}
\usepackage{fontspec}
\setmainfont{Charis SIL}
\newfontfamily\arabicfont[Script=Arabic,Scale=1.2]{Amiri}
\newfontfamily\hebrewfont{Ezra SIL}[Script=Hebrew]
\geometry{margin=1.5cm}
\usepackage{booktabs}
\usepackage{graphicx}
\usepackage{multirow}
\titlespacing*{\section}
  {0pt}                    % left margin
  {3.5ex plus 2ex minus 1.5ex}  % space before (more flexible)
  {2.3ex plus 1ex minus 0.5ex}  % space after (more flexible)

\setlength{\headheight}{14.49998pt}
\addtolength{\topmargin}{-2.49998pt}
\pagestyle{fancy}
\fancyhf{}
\fancyfoot[C]{\thepage}
\fancyhead[R]{\textit{We Are Not in the Abbasid Era}}
\renewcommand{\headrulewidth}{0pt}

\usetikzlibrary{shapes,arrows,decorations.pathmorphing}

% Custom colors
\definecolor{headercolor}{RGB}{70,130,180}
\definecolor{boxcolor}{RGB}{240,248,255}
\definecolor{accentcolor}{RGB}{220,20,60}
\definecolor{tableheader}{RGB}{220,220,220}
\definecolor{dialectcolor}{RGB}{34,139,34}

\begin{document}

\title{\textbf{\Large Arabic Phrase Analyser}\\
\large We Are Not in the Abbasid Era nor in the Umayyad State\\
\normalsize \textarabic{لسنا في العصر العباسي ولسنا في الدولة الأموية}}
\author{Igor Deruga}
\date{}
\maketitle

% ======================== Phrase Display ========================
\begin{tcolorbox}[colback=boxcolor,colframe=headercolor,title=\textbf{Arabic Phrase with Full Diacritics},breakable]
\centering
\Large\textarabic{لَسْنَا فِي الْعَصْرِ الْعَبَّاسِيِّ وَلَسْنَا فِي الدَّوْلَةِ الْأُمَوِيَّةِ}
\end{tcolorbox}

% ======================== Translations ========================
\section{English Translation}
\begin{tcolorbox}[colback=white,colframe=accentcolor,breakable]
\textbf{Literal:} Not-we in the-era the-Abbasid and-not-we in the-state the-Umayyad\\
\textit{[Arabic order retained for direct mapping]}\\[0.5em]
\textbf{Adapted:} We are not in the Abbasid era and we are not in the Umayyad state
\end{tcolorbox}

% ======================== Detailed Word Analysis ========================
\section{Detailed Word Analysis}

\subsection{\textarabic{لَسْنَا} — /lasnaː/}
\begin{tabular}{p{3cm}p{10cm}}
\toprule
\textbf{Translation} & we are not \\
\textbf{Root} & \textarabic{ل-ي-س} (l-j-s) \\
\textbf{Pattern} & Defective verb + pronoun suffix \\
\textbf{Grammar} & Negative verb of being, 1st person plural, with attached pronoun \textarabic{ـنَا} \\
\midrule
\textbf{Examples} & \makecell[l]{\parbox{9.5cm}{
1. \textarabic{لَسْتُ مُتَأَكِّدًا} — I am not sure /lastu mutaʔakkidan/\\
2. \textarabic{لَيْسَ هُنَاكَ مُشْكِلَةٌ} — There is no problem /lajsa hunaːka muʃkilatun/\\
3. \textarabic{لَسْنَا وَحْدَنَا} — We are not alone /lasnaː waħdanaː/
}} \\
\midrule
\textbf{Function} & Negates nominal sentences; acts as auxiliary verb requiring predicate in accusative \\
\textbf{Etymology} & From root meaning "to not be"; unique negative copula in Arabic \\
\bottomrule
\end{tabular}

\subsubsection*{Complete Conjugation of \textarabic{لَيْسَ}}
\par{\large \textbf{Negative Copula - Present Tense Only}}

\begin{longtable}{|>{\raggedright}p{4cm}|p{4cm}|p{5cm}|}
\hline
\textbf{Person} & \textbf{Form} & \textbf{Translation} \\
\hline
\textbf{1st person singular} & \textarabic{لَسْتُ} /lastu/ & I am not \\
\hline
\textbf{1st person plural} & \textarabic{لَسْنَا} /lasnaː/ & we are not \\
\hline
\textbf{2nd person masc. sing.} & \textarabic{لَسْتَ} /lasta/ & you (m.) are not \\
\hline
\textbf{2nd person fem. sing.} & \textarabic{لَسْتِ} /lasti/ & you (f.) are not \\
\hline
\textbf{2nd person dual} & \textarabic{لَسْتُمَا} /lastumaː/ & you (dual) are not \\
\hline
\textbf{2nd person masc. plural} & \textarabic{لَسْتُمْ} /lastum/ & you (m. pl.) are not \\
\hline
\textbf{2nd person fem. plural} & \textarabic{لَسْتُنَّ} /lastunna/ & you (f. pl.) are not \\
\hline
\textbf{3rd person masc. sing.} & \textarabic{لَيْسَ} /lajsa/ & he is not \\
\hline
\textbf{3rd person fem. sing.} & \textarabic{لَيْسَتْ} /lajsat/ & she is not \\
\hline
\textbf{3rd person dual} & \textarabic{لَيْسَا} /lajsaː/ & they (dual) are not \\
\hline
\textbf{3rd person masc. plural} & \textarabic{لَيْسُوا} /lajsuː/ & they (m.) are not \\
\hline
\textbf{3rd person fem. plural} & \textarabic{لَسْنَ} /lasna/ & they (f.) are not \\
\hline
\end{longtable}

\subsubsection*{Grammatical Properties of \textarabic{لَيْسَ}}
\begin{itemize}
  \item \textbf{Category:} Defective verb (\textarabic{فِعْل نَاقِص}) — incomplete verb
  \item \textbf{Function:} Negates nominal sentences only
  \item \textbf{Tense:} Has only present tense form (no past or future conjugations)
  \item \textbf{Case government:}
    \begin{itemize}
      \item Subject (\textarabic{اسم ليس}) — nominative case
      \item Predicate (\textarabic{خبر ليس}) — accusative case
    \end{itemize}
  \item \textbf{Sister verbs:} Part of \textarabic{كَانَ وَأَخَوَاتُهَا} (kāna and its sisters)
  \item \textbf{Example structure:} \textarabic{لَيْسَ الطَّالِبُ حَاضِرًا} — The student is not present
    \begin{itemize}
      \item \textarabic{الطَّالِبُ} (nominative) — subject
      \item \textarabic{حَاضِرًا} (accusative) — predicate
    \end{itemize}
\end{itemize}

\subsection{\textarabic{فِي} — /fiː/}
\begin{tabular}{p{3cm}p{10cm}}
\toprule
\textbf{Translation} & in, at \\
\textbf{Root} & Preposition (not from triliteral root) \\
\textbf{Pattern} & Independent preposition \\
\textbf{Grammar} & Preposition indicating location (literal or metaphorical) \\
\midrule
\textbf{Examples} & \makecell[l]{\parbox{9.5cm}{
1. \textarabic{فِي الْبَيْتِ} — In the house /fiː lbajti/\\
2. \textarabic{فِي الصَّبَاحِ} — In the morning /fiː sˤsˤabaːħi/\\
3. \textarabic{فِي رَأْيِي} — In my opinion /fiː raʔjiː/
}} \\
\midrule
\textbf{Function} & Governs genitive case on following noun \\
\textbf{Etymology} & From Proto-Semitic *fī; Hebrew \texthebrew{בְּ} (be-) "in" \\
\bottomrule
\end{tabular}

\subsubsection*{Usage Nuances}
\begin{itemize}
\item \textbf{Literal location:} \textarabic{فِي الْمَدْرَسَةِ} (in the school)
\item \textbf{Temporal:} \textarabic{فِي الصَّيْفِ} (in summer)
\item \textbf{Abstract/metaphorical:} \textarabic{فِي الْحَقِيقَةِ} (in reality)
\item \textbf{State/condition:} \textarabic{فِي حَالَةٍ جَيِّدَةٍ} (in good condition)
\item \textbf{With eras/periods:} \textarabic{فِي الْعَصْرِ الْحَدِيثِ} (in the modern era)
\end{itemize}

\subsection{\textarabic{الْعَصْرِ} — /alʕasˤri/}
\begin{tabular}{p{3cm}p{10cm}}
\toprule
\textbf{Translation} & the era, the age, the period \\
\textbf{Root} & \textarabic{ع-ص-ر} (ʕ-sˤ-r) \\
\textbf{Pattern} & \textarabic{فَعْل} (faʕl) \\
\textbf{Grammar} & Masculine noun, genitive case (after preposition \textarabic{فِي}), definite \\
\midrule
\textbf{Examples} & \makecell[l]{\parbox{9.5cm}{
1. \textarabic{الْعَصْرُ الذَّهَبِيُّ} — The golden age /alʕasˤru ððahabijju/\\
2. \textarabic{عَصْرُ النَّهْضَةِ} — The Renaissance era /ʕasˤru nnahðˤati/\\
3. \textarabic{فِي هَذَا الْعَصْرِ} — In this age /fiː haːðaː lʕasˤri/
}} \\
\midrule
\textbf{Synonyms} & \textarabic{زَمَن} (time), \textarabic{عَهْد} (epoch), \textarabic{دَهْر} (age), \textarabic{حِقْبَة} (period) \\
\textbf{Etymology} & From root meaning "to press, squeeze"; extended to mean "time period" \\
\bottomrule
\end{tabular}

\subsubsection*{Declension}
\begin{tabular}{|c|c|c|}
\hline
\textbf{Case} & \textbf{Indefinite} & \textbf{Definite} \\
\hline
Nominative & \textarabic{عَصْرٌ} /ʕasˤrun/ & \textarabic{الْعَصْرُ} /alʕasˤru/ \\
\hline
Accusative & \textarabic{عَصْرًا} /ʕasˤran/ & \textarabic{الْعَصْرَ} /alʕasˤra/ \\
\hline
Genitive & \textarabic{عَصْرٍ} /ʕasˤrin/ & \textarabic{الْعَصْرِ} /alʕasˤri/ \\
\hline
\end{tabular}

\subsection{\textarabic{الْعَبَّاسِيِّ} — /alʕabbaːsijji/}
\begin{tabular}{p{3cm}p{10cm}}
\toprule
\textbf{Translation} & the Abbasid \\
\textbf{Root} & \textarabic{ع-ب-س} (ʕ-b-s) from proper name \textarabic{عَبَّاس} \\
\textbf{Pattern} & \textarabic{الْفَعَّالِيّ} (al-faʕʕaːlijj) — definite nisba adjective \\
\textbf{Grammar} & Nisba adjective, masculine singular, genitive case (agreeing with \textarabic{الْعَصْرِ}), definite \\
\midrule
\textbf{Examples} & \makecell[l]{\parbox{9.5cm}{
1. \textarabic{الْخَلِيفَةُ الْعَبَّاسِيُّ} — The Abbasid caliph /alxaliːfatu lʕabbaːsijju/\\
2. \textarabic{الْعَصْرُ الْعَبَّاسِيُّ الْأَوَّلُ} — The first Abbasid era /alʕasˤru lʕabbaːsijju l-ʔawwalu/\\
3. \textarabic{الْحَضَارَةُ الْعَبَّاسِيَّةُ} — The Abbasid civilization /alħadˤaːratu lʕabbaːsijjatu/
}} \\
\midrule
\textbf{Historical Context} & Abbasid Caliphate (750-1258 CE), Islamic Golden Age, capital Baghdad \\
\textbf{Etymology} & From \textarabic{عَبَّاس} (al-ʕAbbās), uncle of Prophet Muhammad \\
\bottomrule
\end{tabular}

\subsubsection*{Declension}
\begin{tabular}{|c|c|c|}
\hline
\textbf{Case} & \textbf{Indefinite} & \textbf{Definite} \\
\hline
Nominative & \textarabic{عَبَّاسِيٌّ} /ʕabbaːsijjun/ & \textarabic{الْعَبَّاسِيُّ} /alʕabbaːsijju/ \\
\hline
Accusative & \textarabic{عَبَّاسِيًّا} /ʕabbaːsijjan/ & \textarabic{الْعَبَّاسِيَّ} /alʕabbaːsijja/ \\
\hline
Genitive & \textarabic{عَبَّاسِيٍّ} /ʕabbaːsijjin/ & \textarabic{الْعَبَّاسِيِّ} /alʕabbaːsijji/ \\
\hline
\end{tabular}

\subsection{\textarabic{وَ} — /wa/}
\begin{tabular}{p{3cm}p{10cm}}
\toprule
\textbf{Translation} & and \\
\textbf{Root} & Coordinating conjunction (not from triliteral root) \\
\textbf{Pattern} & Prefixed conjunction \\
\textbf{Grammar} & Coordinating conjunction connecting two clauses \\
\midrule
\textbf{Examples} & \makecell[l]{\parbox{9.5cm}{
1. \textarabic{مُحَمَّدٌ وَعَلِيٌّ} — Muhammad and Ali /muħammadun wa ʕalijjun/\\
2. \textarabic{قَرَأْتُ وَكَتَبْتُ} — I read and wrote /qaraʔtu wa katabtu/\\
3. \textarabic{الشَّمْسُ وَالْقَمَرُ} — The sun and the moon /aʃʃamsu walqamaru/
}} \\
\midrule
\textbf{Function} & Connects words, phrases, or clauses; most common conjunction in Arabic \\
\textbf{Etymology} & From Proto-Semitic *wa; Hebrew \texthebrew{וְ} (ve-) "and" \\
\bottomrule
\end{tabular}

\subsubsection*{Phonological Variations}
\begin{itemize}
\item Before consonants: \textarabic{وَ} /wa/ — \textarabic{وَمُحَمَّد} /wa muħammad/
\item Before \textarabic{ا}: becomes \textarabic{وَ} — \textarabic{وَالْكِتَاب} /walkitaːb/
\item In connected speech: may reduce to /u/ or /i/ depending on context
\end{itemize}

\subsection{\textarabic{الدَّوْلَةِ} — /ddawlati/}
\begin{tabular}{p{3cm}p{10cm}}
\toprule
\textbf{Translation} & the state, the dynasty, the empire \\
\textbf{Root} & \textarabic{د-و-ل} (d-w-l) \\
\textbf{Pattern} & \textarabic{فَعْلَة} (faʕlah) \\
\textbf{Grammar} & Feminine noun, genitive case (after preposition \textarabic{فِي}), definite \\
\midrule
\textbf{Examples} & \makecell[l]{\parbox{9.5cm}{
1. \textarabic{الدَّوْلَةُ الْعُثْمَانِيَّةُ} — The Ottoman state /addawlatu lʕuθmaːnijjatu/\\
2. \textarabic{رَئِيسُ الدَّوْلَةِ} — Head of state /raʔiːsu ddawlati/\\
3. \textarabic{دَوْلَةٌ عَرَبِيَّةٌ} — An Arab state /dawlatun ʕarabijjatun/
}} \\
\midrule
\textbf{Synonyms} & \textarabic{حُكُومَة} (government), \textarabic{مَمْلَكَة} (kingdom), \textarabic{سُلْطَة} (authority) \\
\textbf{Etymology} & From root meaning "to alternate, rotate"; extended to "ruling power" \\
\bottomrule
\end{tabular}

\subsubsection*{Declension}
\begin{tabular}{|c|c|c|}
\hline
\textbf{Case} & \textbf{Indefinite} & \textbf{Definite} \\
\hline
Nominative & \textarabic{دَوْلَةٌ} /dawlatun/ & \textarabic{الدَّوْلَةُ} /addawlatu/ \\
\hline
Accusative & \textarabic{دَوْلَةً} /dawlatan/ & \textarabic{الدَّوْلَةَ} /addawlata/ \\
\hline
Genitive & \textarabic{دَوْلَةٍ} /dawlatin/ & \textarabic{الدَّوْلَةِ} /addawlati/ \\
\hline
\end{tabular}

\subsection{\textarabic{الْأُمَوِيَّةِ} — /l-ʔumawijjati/}
\begin{tabular}{p{3cm}p{10cm}}
\toprule
\textbf{Translation} & the Umayyad \\
\textbf{Root} & \textarabic{أ-م-ي} from proper name \textarabic{أُمَيَّة} \\
\textbf{Pattern} & \textarabic{الْفُعَلِيَّة} (al-fuʕalijjah) — definite feminine nisba adjective \\
\textbf{Grammar} & Nisba adjective, feminine singular, genitive case (agreeing with \textarabic{الدَّوْلَةِ}), definite \\
\midrule
\textbf{Examples} & \makecell[l]{\parbox{9.5cm}{
1. \textarabic{الْخِلَافَةُ الْأُمَوِيَّةُ} — The Umayyad Caliphate /alxilaːfatu l-ʔumawijjatu/\\
2. \textarabic{الدَّوْلَةُ الْأُمَوِيَّةُ} — The Umayyad state /addawlatu l-ʔumawijjatu/\\
3. \textarabic{الْمَسْجِدُ الْأُمَوِيُّ} — The Umayyad Mosque /almasdʒidu l-ʔumawijju/
}} \\
\midrule
\textbf{Historical Context} & Umayyad Caliphate (661-750 CE), first major Islamic dynasty, capital Damascus \\
\textbf{Etymology} & From \textarabic{أُمَيَّة} (Umayya), ancestor of Umayyad clan \\
\bottomrule
\end{tabular}

\subsubsection*{Declension}
\begin{tabular}{|c|c|c|}
\hline
\textbf{Case} & \textbf{Indefinite} & \textbf{Definite} \\
\hline
Nominative & \textarabic{أُمَوِيَّةٌ} /ʔumawijjatun/ & \textarabic{الْأُمَوِيَّةُ} /al-ʔumawijjatu/ \\
\hline
Accusative & \textarabic{أُمَوِيَّةً} /ʔumawijjatan/ & \textarabic{الْأُمَوِيَّةَ} /al-ʔumawijjata/ \\
\hline
Genitive & \textarabic{أُمَوِيَّةٍ} /ʔumawijjatin/ & \textarabic{الْأُمَوِيَّةِ} /al-ʔumawijjati/ \\
\hline
\end{tabular}

% ======================== Phrase Analysis ========================
\section{Phrase Analysis}

\begin{tcolorbox}[colback=boxcolor,colframe=headercolor,breakable]
\textbf{Grammatical Structure:}\\
Negative copula verb (1st pl.) + prepositional phrase (definite noun + definite adjective) + coordinating conjunction + negative copula verb (1st pl.) + prepositional phrase (definite noun + definite adjective)

\vspace{0.5em}
\textbf{Sentence Type:} Two coordinated negative nominal sentences using \textarabic{لَيْسَ}

\vspace{0.5em}
\textbf{Key Grammar Points:}
\begin{itemize}
\item \textbf{Negative copula:} \textarabic{لَسْنَا} (we are not) — defective verb negating state of being
\item \textbf{Parallel structure:} Two identical sentence structures connected by \textarabic{وَ}
\item \textbf{Prepositional predicates:} Both \textarabic{فِي الْعَصْرِ الْعَبَّاسِيِّ} and \textarabic{فِي الدَّوْلَةِ الْأُمَوِيَّةِ} function as predicates
\item \textbf{Case government by \textarabic{لَيْسَ}:}
  \begin{itemize}
    \item Subject (implicit \textarabic{نَحْنُ}) would be nominative
    \item Predicate (prepositional phrases) treated as manṣūb (accusative) position
  \end{itemize}
\item \textbf{Definiteness:} Both alternatives are fully definite (contrast with previous question phrase)
\item \textbf{Adjective agreement:}
  \begin{itemize}
    \item \textarabic{الْعَبَّاسِيِّ} agrees with \textarabic{الْعَصْرِ} (masc., sing., gen., def.)
    \item \textarabic{الْأُمَوِيَّةِ} agrees with \textarabic{الدَّوْلَةِ} (fem., sing., gen., def.)
  \end{itemize}
\item \textbf{Rhetorical function:} Declarative denial — assertive statement rejecting anachronistic characterization
\item \textbf{Emphasis through repetition:} Repeating \textarabic{لَسْنَا فِي} emphasizes rejection of both alternatives
\end{itemize}
\end{tcolorbox}

\clearpage
\section{Syntactic Analysis}

\begin{tcolorbox}[colback=white,colframe=accentcolor,breakable]
\textbf{Phrase Structure Tree:}

\begin{center}
\begin{tikzpicture}[
  level distance=1.2cm,
  level 1/.style={sibling distance=6cm},
  level 2/.style={sibling distance=3cm},
  level 3/.style={sibling distance=2.5cm},
  every node/.style={rectangle, draw, rounded corners, minimum width=2.2cm, align=center, font=\small}
]

\node {Compound Sentence}
  child {node {Negative\\Sentence 1}
    child {node {\textarabic{لَسْنَا}\\we are not}
      child {node {Implicit\\Subject\\\textarabic{نَحْنُ}}}
    }
    child {node {Predicate}
      child {node {\textarabic{فِي الْعَصْرِ}\\in the era}}
      child {node {\textarabic{الْعَبَّاسِيِّ}\\the Abbasid}}
    }
  }
  child {node {\textarabic{وَ}\\and}}
  child {node {Negative\\Sentence 2}
    child {node {\textarabic{لَسْنَا}\\we are not}
      child {node {Implicit\\Subject\\\textarabic{نَحْنُ}}}
    }
    child {node {Predicate}
      child {node {\textarabic{فِي الدَّوْلَةِ}\\in the state}}
      child {node {\textarabic{الْأُمَوِيَّةِ}\\the Umayyad}}
    }
  };

\end{tikzpicture}
\end{center}

\vspace{1em}
\textbf{Syntactic Relationships:}
\begin{enumerate}
\item \textbf{Main structure:} Two coordinated negative nominal sentences
\item \textbf{Each sentence structure:}
  \begin{itemize}
    \item \textbf{Negative verb:} \textarabic{لَسْنَا} (ism laysa implied: \textarabic{نَحْنُ})
    \item \textbf{Predicate:} Prepositional phrase (khabar laysa)
  \end{itemize}
\item \textbf{Coordination:} \textarabic{وَ} connects two parallel negative statements
\item \textbf{Parallel structure emphasis:} Repetition of \textarabic{لَسْنَا فِي} creates rhythmic emphasis
\item \textbf{Adjective position:} Both nisba adjectives follow their nouns as attributes
\item \textbf{Implicit vs. explicit:} Subject pronoun incorporated in verb, not stated separately
\end{enumerate}

\vspace{1em}
\textbf{Box Diagram:}

\begin{center}
\begin{tikzpicture}[
  box/.style={rectangle, draw, minimum width=2.5cm, minimum height=0.8cm, align=center},
  arrow/.style={->, >=stealth, thick}
]

% First clause
\node[box] (lasna1) at (0,0) {\textarabic{لَسْنَا}\\we are not};
\node[box] (fi1) at (3,0) {\textarabic{فِي}\\in};
\node[box] (3asr) at (5.5,0) {\textarabic{الْعَصْرِ}\\the era};
\node[box] (3abbasi) at (5.5,-2) {\textarabic{الْعَبَّاسِيِّ}\\the Abbasid};

% Conjunction
\node[box, fill=yellow!20] (wa) at (8,-1) {\textarabic{وَ}\\and};

% Second clause
\node[box] (lasna2) at (11,0) {\textarabic{لَسْنَا}\\we are not};
\node[box] (fi2) at (14,0) {\textarabic{فِي}\\in};
\node[box] (dawla) at (16.5,0) {\textarabic{الدَّوْلَةِ}\\the state};
\node[box] (umawi) at (16.5,-2) {\textarabic{الْأُمَوِيَّةِ}\\the Umayyad};

% Arrows
\draw[arrow] (lasna1) -- (fi1);
\draw[arrow] (fi1) -- (3asr);
\draw[arrow] (3asr) -- (3abbasi);
\draw[arrow] (3abbasi) -- (wa);
\draw[arrow] (wa) -- (lasna2);
\draw[arrow] (lasna2) -- (fi2);
\draw[arrow] (fi2) -- (dawla);
\draw[arrow] (dawla) -- (umawi);

\end{tikzpicture}
\end{center}
\end{tcolorbox}

% ======================== Similar Phrases for Practice ========================
\section{Similar Phrases for Practice}

\begin{enumerate}
\item \textarabic{لَسْنَا فِي الْقُرُونِ الْوُسْطَى وَلَسْنَا فِي عَصْرِ الْجَاهِلِيَّةِ}\\
We are not in the Middle Ages and we are not in the era of pre-Islamic ignorance /lasnaː fiː lquruːni lwustˤaː wa lasnaː fiː ʕasˤri ldʒaːhillijjati/

\item \textarabic{لَسْنَا فِي الْخِلَافَةِ الْعُثْمَانِيَّةِ وَلَسْنَا فِي الدَّوْلَةِ الْفَاطِمِيَّةِ}\\
We are not in the Ottoman Caliphate and we are not in the Fatimid state /lasnaː fiː lxilaːfati lʕuθmaːnijjati wa lasnaː fiː ddawlati lfaːtˤimijjati/

\item \textarabic{لَسْتُ فِي الْمَاضِي وَلَسْتُ فِي الْمُسْتَقْبَلِ، أَنَا فِي الْحَاضِرِ}\\
I am not in the past and I am not in the future, I am in the present /lastu fiː lmaːdˤiː wa lastu fiː lmustaqbali, ʔanaː fiː lħaːdˤiri/

\item \textarabic{لَيْسُوا فِي عَهْدِ الْمَمَالِيكِ وَلَيْسُوا فِي عَهْدِ الْأَيُّوبِيِّينَ}\\
They are not in the Mamluk epoch and they are not in the Ayyubid epoch /lajsuː fiː ʕahdi lmamaːliːki wa lajsuː fiː ʕahdi l-ʔajjuːbijjiːna/

\item \textarabic{لَسْنَا فِي زَمَنِ الْاِسْتِعْمَارِ وَلَسْنَا فِي زَمَنِ الْعُبُودِيَّةِ}\\
We are not in the time of colonialism and we are not in the time of slavery /lasnaː fiː zamani l-istiʕmaːri wa lasnaː fiː zamani lʕubuːdijjati/

\item \textarabic{لَسْتَ فِي بَغْدَادَ الْعَبَّاسِيَّةِ وَلَسْتَ فِي دِمَشْقَ الْأُمَوِيَّةِ}\\
You are not in Abbasid Baghdad and you are not in Umayyad Damascus /lasta fiː baɣdaːda lʕabbaːsijjati wa lasta fiː dimaʃqa l-ʔumawijjati/
\end{enumerate}

% ======================== Levantine Dialect Version ========================
\section{Levantine (Shami) Arabic Dialect}

\begin{tcolorbox}[colback=white,colframe=dialectcolor,title=\textbf{Levantine Version},breakable]
\textarabic{مِش نِحْنَا بِالْعَصْر الْعَبَّاسِي وَمِش نِحْنَا بِالدَّوْلِة الْأُمَوِيِّة}\\[0.5em]
\textbf{Phonetic:} /miʃ niħna bilʕasˤar elʕabbaːsiː wmiʃ niħna biddawle l-ʔumawijje/\\[0.5em]
\textbf{Translation:} We are not in the Abbasid era and we are not in the Umayyad state
\end{tcolorbox}

\textbf{Key Dialectal Changes:}
\begin{itemize}
\item \textarabic{لَسْنَا} → \textarabic{مِش نِحْنَا} /miʃ niħna/ — Dialectal negation with explicit pronoun
\item Alternative: \textarabic{مَا نِحْنَا} /maː niħna/ — Another dialectal negation pattern
\item \textarabic{فِي} → \textarabic{بِ} /bi/ — Preposition shift (common in Levantine)
\item \textarabic{الْعَصْرِ} → \textarabic{الْعَصْر} /elʕasˤar/ — Loss of case endings
\item \textarabic{الْعَبَّاسِيِّ} → \textarabic{الْعَبَّاسِي} /elʕabbaːsiː/ — Loss of case endings
\item \textarabic{وَلَسْنَا} → \textarabic{وَمِش} /wmiʃ/ — Shortened conjunction + negation
\item \textarabic{الدَّوْلَةِ} → \textarabic{الدَّوْلِة} /eddawle/ — Dialectal pronunciation
\item \textarabic{الْأُمَوِيَّةِ} → \textarabic{الْأُمَوِيِّة} /l-ʔumawijje/ — Simplified ending
\end{itemize}

\subsection{Alternative Dialectal Versions}

\textbf{Egyptian Arabic:}\\
\textarabic{إِحْنَا مِش فِي الْعَصْر الْعَبَّاسِي وَمِش فِي الدَّوْلَة الْأُمَوِيَّة}\\
/ʔeħna miʃ fi lʕasˤr elʕabbaːsi wmiʃ fi ddawla l-ʔumawejja/

\textbf{Gulf Arabic:}\\
\textarabic{مَا نِحْنَا فِي الْعَصْر الْعَبَّاسِي وَمَا نِحْنَا فِي الدَّوْلَة الْأُمَوِيَّة}\\
/maː niħna fi lʕasˤir elʕabbaːsi wmaː niħna fi ddawla l-ʔumawejja/

\textbf{Moroccan Arabic:}\\
\textarabic{مَا حْنَا شْ فْ الْعَصْر الْعَبَّاسِي وَلَا فْ الدَّوْلَة الْأُمَوِيَّة}\\
/maː ħnaːʃ f lʕasˤr lʕabbaːsi wəla f ddəwla l-ʔumawejja/

\textbf{Iraqi Arabic:}\\
\textarabic{مَاكُو نِحْنَا بِالْعَصْر الْعَبَّاسِي وَلَا بِالدَّوْلَة الْأُمَوِيَّة}\\
/maːku niħna bilʕasˤar elʕabbaːsi wla biddawla l-ʔumawejja/

% ======================== Additional Learning Notes ========================
\section{Additional Learning Notes}

\begin{tcolorbox}[colback=boxcolor,colframe=accentcolor,title=\textbf{Cultural and Rhetorical Context},breakable]

\textbf{Rhetorical Function:}\\
This phrase serves as a declarative rebuttal to the earlier question. It represents:
\begin{itemize}
\item \textbf{Assertion of modernity:} Rejecting anachronistic characterizations
\item \textbf{Political statement:} Denying that current governance resembles medieval systems
\item \textbf{Progressive stance:} Affirming modern identity and rejecting backward-looking critiques
\item \textbf{Defensive response:} Countering accusations of outdated political structures
\end{itemize}

\textbf{Discourse Context:}\\
This phrase typically appears in:
\begin{enumerate}
\item \textbf{Political speeches:} Leaders defending their governance models
\item \textbf{Media commentary:} Journalists discussing modernization efforts
\item \textbf{Academic discourse:} Scholars distinguishing contemporary from historical systems
\item \textbf{Public debates:} Citizens asserting modern political consciousness
\item \textbf{Reform discussions:} Activists arguing for progressive change
\end{enumerate}

\textbf{Implied Continuation:}\\
The phrase often continues with assertions about what era/system we ARE in:
\begin{itemize}
\item \textarabic{نَحْنُ فِي عَصْرٍ حَدِيثٍ} — We are in a modern era
\item \textarabic{نَحْنُ فِي الْقَرْنِ الْوَاحِدِ وَالْعِشْرِينَ} — We are in the 21st century
\item \textarabic{نَحْنُ فِي عَصْرِ الدِّيمُقْرَاطِيَّةِ} — We are in the era of democracy
\end{itemize}

\textbf{Contrast with Question Form:}\\
Compare this declarative with the earlier interrogative:
\begin{itemize}
\item Question: \textarabic{هَلْ نَحْنُ بِعَصْرٍ عَبَّاسِيٍّ أَمْ عَصْرِ الدَّوْلَةِ الْأُمَوِيَّةِ؟}
\item Answer/Rebuttal: \textarabic{لَسْنَا فِي الْعَصْرِ الْعَبَّاسِيِّ وَلَسْنَا فِي الدَّوْلَةِ الْأُمَوِيَّةِ}
\item The declarative is more assertive and emphatic through repetition
\end{itemize}

\textbf{Historical Irony:}\\
The phrase gains additional meaning from historical context:
\begin{itemize}
\item Umayyad and Abbasid dynasties represent different political models
\item Both were eventually replaced by modern nation-states
\item The phrase acknowledges this transition while rejecting nostalgia
\item Implicitly recognizes that those systems are incompatible with modernity
\end{itemize}
\end{tcolorbox}

\subsection{Memory Tips}

\begin{enumerate}
\item \textbf{Negation Pattern:} Remember \textarabic{لَيْسَ} conjugates like a verb but only in present tense
  \begin{itemize}
    \item \textarabic{لَسْتُ} (I am not), \textarabic{لَسْنَا} (we are not), \textarabic{لَسْتَ} (you are not)
    \item No past or future forms — use \textarabic{لَمْ يَكُنْ} (was not) for past
  \end{itemize}

\item \textbf{Parallel Structure:} The repetition of \textarabic{لَسْنَا فِي} creates emphasis
  \begin{itemize}
    \item Pattern: [NEG + PREP + NOUN + ADJ] + [CONJ] + [NEG + PREP + NOUN + ADJ]
    \item Repeating the entire structure makes the denial stronger
  \end{itemize}

\item \textbf{Definiteness Change:} Compare with the question:
  \begin{itemize}
    \item Question had: \textarabic{بِعَصْرٍ عَبَّاسِيٍّ} (indefinite — "an Abbasid era")
    \item Answer has: \textarabic{فِي الْعَصْرِ الْعَبَّاسِيِّ} (definite — "the Abbasid era")
    \item Definiteness adds specificity and authority to the denial
  \end{itemize}

\item \textbf{Preposition Switch:}
  \begin{itemize}
    \item Question used \textarabic{بِ} (in/with — more metaphorical)
    \item Answer uses \textarabic{فِي} (in — more literal/locational)
    \item This makes the denial more concrete and grounded
  \end{itemize}

\item \textbf{Verb \textarabic{لَيْسَ} Mnemonic:} Think "Lay-sa" = "Lay (down) the negative" — it lays down a negative state

\item \textbf{Dynasty Contrast:} Remember the key difference:
  \begin{itemize}
    \item Umayyad (661-750) = Damascus = Arab aristocracy
    \item Abbasid (750-1258) = Baghdad = Cosmopolitan/Persian influence
  \end{itemize}
\end{enumerate}

\subsection{Related Vocabulary Families}

\begin{multicols}{2}

\textbf{Negation Words:}\\
\textarabic{لَيْسَ} — is not /lajsa/\\
\textarabic{لَا} — no, not /laː/\\
\textarabic{مَا} — not (verbal negation) /maː/\\
\textarabic{لَمْ} — did not /lam/\\
\textarabic{لَنْ} — will not /lan/\\
\textarabic{غَيْر} — non-, not /ɣajr/\\
\textarabic{بِلَا} — without /bilaː/\\
\textarabic{بِدُونِ} — without /biduːni/\\

\textbf{Time/Era Words:}\\
\textarabic{عَصْر} — era /ʕasˤr/\\
\textarabic{زَمَن} — time /zaman/\\
\textarabic{دَهْر} — age /dahr/\\
\textarabic{عَهْد} — epoch /ʕahd/\\
\textarabic{حِقْبَة} — period /ħiqba/\\
\textarabic{مَرْحَلَة} — stage /marħala/\\
\textarabic{فَتْرَة} — period /fatra/\\
\textarabic{وَقْت} — time /waqt/\\

\textbf{State/Government:}\\
\textarabic{دَوْلَة} — state /dawla/\\
\textarabic{حُكُومَة} — government /ħukuːma/\\
\textarabic{خِلَافَة} — caliphate /xilaːfa/\\
\textarabic{إِمْبَرَاطُورِيَّة} — empire /ʔimbraːtˤuːrijja/\\
\textarabic{مَمْلَكَة} — kingdom /mamlaka/\\
\textarabic{سُلْطَنَة} — sultanate /sultˤana/\\
\textarabic{جُمْهُورِيَّة} — republic /dʒumhuːrijja/\\
\textarabic{نِظَام} — regime /nidˤaːm/\\

\textbf{Historical Adjectives:}\\
\textarabic{قَدِيم} — old, ancient /qadiːm/\\
\textarabic{حَدِيث} — modern /ħadiːθ/\\
\textarabic{مَاضٍ} — past /maːdˤin/\\
\textarabic{حَاضِر} — present /ħaːdˤir/\\
\textarabic{مُسْتَقْبَلِيّ} — future /mustaqbalijj/\\
\textarabic{عَتِيق} — antique /ʕatiːq/\\
\textarabic{مُعَاصِر} — contemporary /muʕaːsˤir/\\
\textarabic{تَقْلِيدِيّ} — traditional /taqliːdijj/\\

\textbf{Conjunctions:}\\
\textarabic{وَ} — and /wa/\\
\textarabic{أَوْ} — or /ʔaw/\\
\textarabic{أَمْ} — or (questions) /ʔam/\\
\textarabic{لَكِنْ} — but /laːkin/\\
\textarabic{بَلْ} — rather /bal/\\
\textarabic{فَ} — so, then /fa/\\
\textarabic{ثُمَّ} — then /θumma/\\
\textarabic{حَتَّى} — until, even /ħattaː/\\

\textbf{Prepositions:}\\
\textarabic{فِي} — in /fiː/\\
\textarabic{بِ} — in, with /bi/\\
\textarabic{مِنْ} — from /min/\\
\textarabic{إِلَى} — to /ʔilaː/\\
\textarabic{عَلَى} — on /ʕalaː/\\
\textarabic{عَنْ} — about, from /ʕan/\\
\textarabic{مَعَ} — with /maʕa/\\
\textarabic{لِ} — for /li/\\
\end{multicols}

\section{Grammatical Deep Dive}

\subsection{The Negative Copula \textarabic{لَيْسَ}}

\begin{tcolorbox}[colback=boxcolor,colframe=headercolor,breakable]
\textbf{Special Status of \textarabic{لَيْسَ}:}

\textarabic{لَيْسَ} is unique in Arabic grammar:
\begin{itemize}
\item \textbf{Classification:} One of \textarabic{كَانَ وَأَخَوَاتُهَا} (kāna and its sisters)
\item \textbf{Function:} Negates nominal (equational) sentences
\item \textbf{Tense limitation:} Only present tense (no past or future)
\item \textbf{Case government:} Takes nominative subject and accusative predicate
\end{itemize}

\textbf{Contrast with Other Negation:}

\begin{tabular}{|l|l|l|}
\hline
\textbf{Particle} & \textbf{Usage} & \textbf{Example} \\
\hline
\textarabic{لَيْسَ} & Nominal sentences (present) & \textarabic{لَسْتُ طَالِبًا} (I am not a student) \\
\hline
\textarabic{لَا} & Categorical negation & \textarabic{لَا إِلَهَ إِلَّا اللهُ} (No god but God) \\
\hline
\textarabic{مَا} & Verbal negation (past) & \textarabic{مَا قَرَأْتُ} (I did not read) \\
\hline
\textarabic{لَمْ} & Verbal negation (jussive) & \textarabic{لَمْ أَقْرَأْ} (I did not read) \\
\hline
\textarabic{لَنْ} & Future negation (subjunctive) & \textarabic{لَنْ أَقْرَأَ} (I will not read) \\
\hline
\end{tabular}

\vspace{0.5em}
\textbf{Case Government Example:}

\textarabic{لَيْسَ الطَّالِبُ حَاضِرًا}\\
/lajsa tˤtˤaːlibu ħaːdˤiran/\\
The student is not present

\begin{itemize}
\item \textarabic{لَيْسَ} — negative copula
\item \textarabic{الطَّالِبُ} — subject (ism laysa), nominative
\item \textarabic{حَاضِرًا} — predicate (khabar laysa), accusative
\end{itemize}

\textbf{Sisters of \textarabic{كَانَ}:}

These verbs all govern nominative subject + accusative predicate:
\begin{itemize}
\item \textarabic{كَانَ} — was, is
\item \textarabic{أَصْبَحَ} — became (morning)
\item \textarabic{أَضْحَى} — became (forenoon)
\item \textarabic{أَمْسَى} — became (evening)
\item \textarabic{ظَلَّ} — remained
\item \textarabic{بَاتَ} — spent the night
\item \textarabic{صَارَ} — became
\item \textarabic{لَيْسَ} — is not
\item \textarabic{مَا زَالَ} — still is
\item \textarabic{مَا بَرِحَ} — has not ceased
\item \textarabic{مَا فَتِئَ} — has not stopped
\item \textarabic{مَا انْفَكَّ} — has not separated from
\item \textarabic{مَا دَامَ} — as long as
\end{itemize}
\end{tcolorbox}

\subsection{Prepositional Predicates}

\begin{tcolorbox}[colback=white,colframe=accentcolor,breakable]
\textbf{Prepositional Phrases as Predicates:}

In our phrase: \textarabic{لَسْنَا فِي الْعَصْرِ الْعَبَّاسِيِّ}

The prepositional phrase \textarabic{فِي الْعَصْرِ الْعَبَّاسِيِّ} functions as the predicate (khabar).

\textbf{Structure Analysis:}
\begin{itemize}
\item \textbf{Subject:} Implicit \textarabic{نَحْنُ} (we) — incorporated in \textarabic{لَسْنَا}
\item \textbf{Predicate:} \textarabic{فِي الْعَصْرِ الْعَبَّاسِيِّ} — prepositional phrase
\item \textbf{Status:} The entire PP is treated as manṣūb (accusative position)
\end{itemize}

\textbf{Why Prepositional Phrases Work as Predicates:}

Arabic allows PP predicates because they:
\begin{enumerate}
\item Provide location information about the subject
\item Complete the meaning of the sentence
\item Function semantically like adjectives or nouns
\end{enumerate}

\textbf{Examples of PP Predicates:}

\begin{tabular}{|l|l|}
\hline
\textbf{Arabic} & \textbf{English} \\
\hline
\textarabic{الْكِتَابُ عَلَى الطَّاوِلَةِ} & The book is on the table \\
\hline
\textarabic{نَحْنُ فِي الْمَدْرَسَةِ} & We are at school \\
\hline
\textarabic{هُوَ مِنْ مِصْرَ} & He is from Egypt \\
\hline
\textarabic{الْقَلَمُ بِيَدِي} & The pen is in my hand \\
\hline
\end{tabular}

\textbf{With \textarabic{لَيْسَ}:}

\begin{tabular}{|l|l|}
\hline
\textbf{Arabic} & \textbf{English} \\
\hline
\textarabic{لَيْسَ الْكِتَابُ عَلَى الطَّاوِلَةِ} & The book is not on the table \\
\hline
\textarabic{لَسْنَا فِي الْمَدْرَسَةِ} & We are not at school \\
\hline
\textarabic{لَيْسَ هُوَ مِنْ مِصْرَ} & He is not from Egypt \\
\hline
\end{tabular}
\end{tcolorbox}

\subsection{Coordination and Parallel Structure}

\begin{tcolorbox}[colback=boxcolor,colframe=headercolor,breakable]
\textbf{Parallel Structure in Our Phrase:}

\textarabic{لَسْنَا فِي الْعَصْرِ الْعَبَّاسِيِّ وَلَسْنَا فِي الدَّوْلَةِ الْأُمَوِيَّةِ}

\textbf{Pattern Analysis:}
\begin{center}
\begin{tabular}{|c|c|c|}
\hline
\textbf{Element 1} & \textbf{Conjunction} & \textbf{Element 2} \\
\hline
\textarabic{لَسْنَا فِي X} & \textarabic{وَ} & \textarabic{لَسْنَا فِي Y} \\
\hline
\end{tabular}
\end{center}

\textbf{Rhetorical Effect of Repetition:}
\begin{itemize}
\item \textbf{Emphasis:} Repeating \textarabic{لَسْنَا} strengthens the denial
\item \textbf{Rhythm:} Creates parallel cadence for memorable impact
\item \textbf{Completeness:} Each denial stands independently
\item \textbf{Clarity:} No ambiguity about what is being negated
\end{itemize}

\textbf{Alternative: Elliptical Structure}

Could be simplified to:
\textarabic{لَسْنَا فِي الْعَصْرِ الْعَبَّاسِيِّ وَلَا فِي الدَّوْلَةِ الْأُمَوِيَّةِ}

\begin{itemize}
\item Omits second \textarabic{لَسْنَا}
\item Uses \textarabic{وَلَا} (and not) instead
\item More concise but less emphatic
\item Still grammatically correct
\end{itemize}

\textbf{Coordination Patterns:}

\begin{tabular}{|l|p{8cm}|}
\hline
\textbf{Type} & \textbf{Example} \\
\hline
Full repetition & \textarabic{لَسْنَا... وَلَسْنَا...} (most emphatic) \\
\hline
Partial ellipsis & \textarabic{لَسْنَا... وَلَا...} (less emphatic) \\
\hline
Shared verb & \textarabic{لَسْنَا... أَوْ...} (for alternatives) \\
\hline
\end{tabular}
\end{tcolorbox}

\section{Practice Exercises}

\subsection{Translation Practice}

Translate into Arabic using \textarabic{لَيْسَ} structure:

\begin{enumerate}
\item I am not in the past and I am not in the future
\vspace{2em}

\item They are not in the Ottoman era and they are not in the Fatimid state
\vspace{2em}

\item You (fem. sing.) are not in Damascus and you are not in Baghdad
\vspace{2em}

\item We are not children and we are not ignorant
\vspace{2em}
\end{enumerate}

\subsection{Grammatical Transformation}

Transform the original phrase:

\begin{enumerate}
\item Change subject to 3rd person masculine plural (\textarabic{هُمْ}):
\vspace{2em}

\item Change to affirmative using \textarabic{كَانَ} (past tense):
\vspace{2em}

\item Use elliptical structure (omit second \textarabic{لَسْنَا}):
\vspace{2em}

\item Add temporal marker \textarabic{الْآنَ} (now):
\vspace{2em}
\end{enumerate}

\subsection{Case Identification}

Identify the grammatical case and explain:

\begin{enumerate}
\item \textarabic{لَسْنَا فِي الْ\underline{عَصْرِ} الْعَبَّاسِيِّ}\\
Case: \underline{\hspace{8cm}}\\
Reason: \underline{\hspace{8cm}}

\item \textarabic{لَسْنَا فِي الْعَصْرِ الْ\underline{عَبَّاسِيِّ}}\\
Case: \underline{\hspace{8cm}}\\
Reason: \underline{\hspace{8cm}}

\item \textarabic{لَسْنَا فِي الْ\underline{دَوْلَةِ} الْأُمَوِيَّةِ}\\
Case: \underline{\hspace{8cm}}\\
Reason: \underline{\hspace{8cm}}
\end{enumerate}

\subsection{Negation Conversion}

Convert these affirmative sentences to negative using \textarabic{لَيْسَ}:

\begin{enumerate}
\item \textarabic{نَحْنُ فِي الْقَرْنِ الْعِشْرِينَ}\\
Negative: \underline{\hspace{10cm}}

\item \textarabic{هُوَ طَالِبٌ}\\
Negative: \underline{\hspace{10cm}}

\item \textarabic{هِيَ فِي الْمَدْرَسَةِ}\\
Negative: \underline{\hspace{10cm}}

\item \textarabic{أَنْتُمْ فِي بَغْدَادَ}\\
Negative: \underline{\hspace{10cm}}
\end{enumerate}

\section{Answers to Practice Exercises}

\subsection{Translation Practice — Suggested Answers}

\begin{enumerate}
\item \textbf{I am not in the past and I am not in the future}\\
\textarabic{لَسْتُ فِي الْمَاضِي وَلَسْتُ فِي الْمُسْتَقْبَلِ}\\
/lastu fiː lmaːdˤiː wa lastu fiː lmustaqbali/

\item \textbf{They are not in the Ottoman era and they are not in the Fatimid state}\\
\textarabic{لَيْسُوا فِي الْعَصْرِ الْعُثْمَانِيِّ وَلَيْسُوا فِي الدَّوْلَةِ الْفَاطِمِيَّةِ}\\
/lajsuː fiː lʕasˤri lʕuθmaːnijji wa lajsuː fiː ddawlati lfaːtˤimijjati/

\item \textbf{You (fem. sing.) are not in Damascus and you are not in Baghdad}\\
\textarabic{لَسْتِ فِي دِمَشْقَ وَلَسْتِ فِي بَغْدَادَ}\\
/lasti fiː dimaʃqa wa lasti fiː baɣdaːda/

\item \textbf{We are not children and we are not ignorant}\\
\textarabic{لَسْنَا أَطْفَالًا وَلَسْنَا جَاهِلِينَ}\\
/lasnaː ʔatˤfaːlan wa lasnaː dʒaːhiliːna/
\end{enumerate}

\subsection{Grammatical Transformation — Answers}

\begin{enumerate}
\item \textbf{3rd person masculine plural:}\\
\textarabic{لَيْسُوا فِي الْعَصْرِ الْعَبَّاسِيِّ وَلَيْسُوا فِي الدَّوْلَةِ الْأُمَوِيَّةِ}\\
/lajsuː fiː lʕasˤri lʕabbaːsijji wa lajsuː fiː ddawlati l-ʔumawijjati/\\
They are not in the Abbasid era and they are not in the Umayyad state

\item \textbf{Affirmative with \textarabic{كَانَ} (past):}\\
\textarabic{كُنَّا فِي الْعَصْرِ الْعَبَّاسِيِّ وَكُنَّا فِي الدَّوْلَةِ الْأُمَوِيَّةِ}\\
/kunnaː fiː lʕasˤri lʕabbaːsijji wa kunnaː fiː ddawlati l-ʔumawijjati/\\
We were in the Abbasid era and we were in the Umayyad state

\item \textbf{Elliptical structure:}\\
\textarabic{لَسْنَا فِي الْعَصْرِ الْعَبَّاسِيِّ وَلَا فِي الدَّوْلَةِ الْأُمَوِيَّةِ}\\
/lasnaː fiː lʕasˤri lʕabbaːsijji wa laː fiː ddawlati l-ʔumawijjati/\\
We are not in the Abbasid era nor in the Umayyad state

\item \textbf{With temporal marker \textarabic{الْآنَ}:}\\
\textarabic{لَسْنَا الْآنَ فِي الْعَصْرِ الْعَبَّاسِيِّ وَلَسْنَا فِي الدَّوْلَةِ الْأُمَوِيَّةِ}\\
/lasnaː l-ʔaːna fiː lʕasˤri lʕabbaːsijji wa lasnaː fiː ddawlati l-ʔumawijjati/\\
We are not now in the Abbasid era and we are not in the Umayyad state
\end{enumerate}

\subsection{Case Identification — Answers}

\begin{enumerate}
\item \textbf{\textarabic{الْعَصْرِ}}: \textbf{Genitive (majrūr)}\\
Reason: Genitive case because it follows the preposition \textarabic{فِي}. All nouns after prepositions must be in the genitive case.

\item \textbf{\textarabic{الْعَبَّاسِيِّ}}: \textbf{Genitive (majrūr)}\\
Reason: Genitive case because it is an adjective modifying \textarabic{الْعَصْرِ}, and adjectives must agree with their nouns in case, gender, number, and definiteness.

\item \textbf{\textarabic{الدَّوْلَةِ}}: \textbf{Genitive (majrūr)}\\
Reason: Genitive case because it follows the preposition \textarabic{فِي}. The feminine singular genitive ending is \textarabic{ـِ} (kasrah).
\end{enumerate}

\subsection{Negation Conversion — Answers}

\begin{enumerate}
\item \textbf{Affirmative:} \textarabic{نَحْنُ فِي الْقَرْنِ الْعِشْرِينَ}\\
\textbf{Negative:} \textarabic{لَسْنَا فِي الْقَرْنِ الْعِشْرِينَ}\\
/lasnaː fiː lqarni lʕiʃriːna/\\
We are not in the twentieth century

\item \textbf{Affirmative:} \textarabic{هُوَ طَالِبٌ}\\
\textbf{Negative:} \textarabic{لَيْسَ هُوَ طَالِبًا} or \textarabic{لَيْسَ طَالِبًا}\\
/lajsa huwa tˤaːliban/ or /lajsa tˤaːliban/\\
He is not a student\\
(Note: Predicate \textarabic{طَالِبًا} becomes accusative)

\item \textbf{Affirmative:} \textarabic{هِيَ فِي الْمَدْرَسَةِ}\\
\textbf{Negative:} \textarabic{لَيْسَتْ هِيَ فِي الْمَدْرَسَةِ} or \textarabic{لَيْسَتْ فِي الْمَدْرَسَةِ}\\
/lajsat hija fiː lmadrasti/ or /lajsat fiː lmadrasti/\\
She is not in the school

\item \textbf{Affirmative:} \textarabic{أَنْتُمْ فِي بَغْدَادَ}\\
\textbf{Negative:} \textarabic{لَسْتُمْ فِي بَغْدَادَ}\\
/lastum fiː baɣdaːda/\\
You (pl.) are not in Baghdad
\end{enumerate}

\section{Extended Context: Discourse Analysis}

\begin{tcolorbox}[colback=boxcolor,colframe=accentcolor,title=\textbf{Question-Answer Dialogue},breakable]

\textbf{Complete Conversational Exchange:}

\textbf{Speaker A (Critic):}\\
\textarabic{هَلْ نَحْنُ بِعَصْرٍ عَبَّاسِيٍّ أَمْ عَصْرِ الدَّوْلَةِ الْأُمَوِيَّةِ؟}\\
Are we in an Abbasid era or the era of the Umayyad state?\\
\textit{[Rhetorical question implying outdated governance]}

\textbf{Speaker B (Defender):}\\
\textarabic{لَسْنَا فِي الْعَصْرِ الْعَبَّاسِيِّ وَلَسْنَا فِي الدَّوْلَةِ الْأُمَوِيَّةِ}\\
We are not in the Abbasid era and we are not in the Umayyad state\\
\textit{[Direct rebuttal asserting modernity]}

\textbf{Possible Continuation:}\\
\textarabic{نَحْنُ فِي الْقَرْنِ الْوَاحِدِ وَالْعِشْرِينَ}\\
We are in the twenty-first century\\
\textit{[Positive assertion of contemporary identity]}

\vspace{0.5em}
\textbf{Rhetorical Analysis:}

\begin{tabular}{|l|p{5cm}|p{5cm}|}
\hline
\textbf{Feature} & \textbf{Question} & \textbf{Answer} \\
\hline
Structure & Interrogative with alternatives & Declarative with parallel negations \\
\hline
Definiteness & Mixed (indef. + def.) & Both definite \\
\hline
Preposition & \textarabic{بِ} (metaphorical) & \textarabic{فِي} (literal) \\
\hline
Tone & Skeptical, accusatory & Defensive, assertive \\
\hline
Function & Challenge authority & Defend modernity \\
\hline
Implication & System is outdated & System is contemporary \\
\hline
\end{tabular}

\vspace{0.5em}
\textbf{Linguistic Strategies:}

\textbf{In the Question:}
\begin{itemize}
\item Uses indefinite for first option (generic criticism)
\item Creates false dilemma (either/or)
\item Implies no modern alternative exists
\item Rhetorical tone suggests answer is "yes"
\end{itemize}

\textbf{In the Answer:}
\begin{itemize}
\item Uses definite articles (specific historical periods)
\item Employs parallel structure for emphasis
\item Categorically rejects both alternatives
\item Creates space for modern identity assertion
\item More literal preposition grounds response in reality
\end{itemize}
\end{tcolorbox}

\section{Comparative Analysis: Question vs. Statement}

\begin{tcolorbox}[colback=white,colframe=accentcolor,breakable]
\textbf{Side-by-Side Comparison:}

\begin{center}
\begin{tabular}{|l|p{6cm}|p{6cm}|}
\hline
\textbf{Feature} & \textbf{Question Form} & \textbf{Statement Form} \\
\hline
\textbf{Sentence Type} & Interrogative & Declarative \\
\hline
\textbf{Opening} & \textarabic{هَلْ} (question particle) & \textarabic{لَسْنَا} (negative verb) \\
\hline
\textbf{Subject} & \textarabic{نَحْنُ} (explicit) & Implicit in verb \\
\hline
\textbf{Preposition} & \textarabic{بِ} & \textarabic{فِي} \\
\hline
\textbf{First Option} & \textarabic{بِعَصْرٍ عَبَّاسِيٍّ} (indef.) & \textarabic{فِي الْعَصْرِ الْعَبَّاسِيِّ} (def.) \\
\hline
\textbf{Connector} & \textarabic{أَمْ} (or) & \textarabic{وَلَسْنَا} (and we are not) \\
\hline
\textbf{Second Option} & \textarabic{عَصْرِ الدَّوْلَةِ...} (def.) & \textarabic{فِي الدَّوْلَةِ...} (def.) \\
\hline
\textbf{Structure} & Alternative question & Parallel negations \\
\hline
\textbf{Verb Repetition} & None (elliptical) & Full (\textarabic{لَسْنَا...وَلَسْنَا}) \\
\hline
\textbf{Tone} & Skeptical, challenging & Assertive, defensive \\
\hline
\textbf{Function} & Criticize & Refute \\
\hline
\textbf{Expectation} & Invites agreement & Demands acceptance \\
\hline
\end{tabular}
\end{center}

\vspace{0.5em}
\textbf{Grammatical Differences:}

\begin{enumerate}
\item \textbf{Question uses \textarabic{هَلْ} + nominal sentence:}
\begin{itemize}
\item \textarabic{هَلْ نَحْنُ بِعَصْرٍ...} — neutral questioning
\item Expects yes/no or alternative answer
\end{itemize}

\item \textbf{Statement uses \textarabic{لَيْسَ} negative verb:}
\begin{itemize}
\item \textarabic{لَسْنَا فِي الْعَصْرِ...} — categorical negation
\item Asserts negative fact, no response expected
\end{itemize}

\item \textbf{Definiteness shift:}
\begin{itemize}
\item Question: \textarabic{عَصْرٍ عَبَّاسِيٍّ} — "an Abbasid era" (any)
\item Statement: \textarabic{الْعَصْرِ الْعَبَّاسِيِّ} — "the Abbasid era" (specific)
\item Definiteness makes denial more concrete and authoritative
\end{itemize}

\item \textbf{Preposition choice:}
\begin{itemize}
\item \textarabic{بِ} — more abstract, metaphorical ("in/with circumstances of")
\item \textarabic{فِي} — more concrete, locational ("located in")
\item Shift to \textarabic{فِي} makes denial more grounded in reality
\end{itemize}
\end{enumerate}
\end{tcolorbox}

\section{Modern Political Discourse Usage}

\begin{tcolorbox}[colback=boxcolor,colframe=headercolor,breakable]
\textbf{Contemporary Context:}

These phrases appear frequently in modern Arab political discourse, particularly in discussions about:

\textbf{1. Governance Models:}
\begin{itemize}
\item Debates about democracy vs. authoritarianism
\item Discussions of hereditary vs. elected leadership
\item Critiques of dynastic political systems
\item Calls for modernization and reform
\end{itemize}

\textbf{2. Historical Consciousness:}
\begin{itemize}
\item Arab societies' relationship with their past
\item Tension between tradition and modernity
\item Islamic heritage vs. secular governance
\item National identity formation
\end{itemize}

\textbf{3. Reform Movements:}
\begin{itemize}
\item Arab Spring rhetoric (2011+)
\item Youth movements demanding change
\item Intellectual debates about progress
\item Civil society advocacy
\end{itemize}

\textbf{Example Contexts:}

\textbf{Political Speech:}\\
\textit{"Critics claim we live in a medieval system, but this is false. We are not in the Abbasid era and we are not in the Umayyad state. We are building a modern nation based on contemporary principles."}

\textbf{Op-Ed Article:}\\
\textit{"The question 'Are we in an Abbasid or Umayyad era?' reflects deep frustration with governance. The answer should be obvious: we are not in either. Yet the question persists because many still feel governed by antiquated systems."}

\textbf{Social Media Debate:}\\
\textit{User 1: "Which medieval caliphate are we living in?"\\
User 2: "We're not in the Abbasid era and not in the Umayyad state. Stop the false comparison."\\
User 3: "Then why does power pass from father to son?"}

\textbf{Academic Analysis:}\\
\textit{"While formally modern nation-states, several Arab countries exhibit governance patterns reminiscent of historical dynasties, prompting the recurring question about which era we truly inhabit."}
\end{tcolorbox}

\section{Cultural Notes: Historical Memory}

\begin{tcolorbox}[colback=white,colframe=accentcolor,breakable]
\textbf{Why These Specific Dynasties?}

The Umayyad and Abbasid periods are referenced because they:

\textbf{1. Historical Significance:}
\begin{itemize}
\item Represent the formative centuries of Islamic civilization
\item Created governmental structures that lasted for centuries
\item Established cultural and religious norms
\item Defined Arabic as language of administration and culture
\end{itemize}

\textbf{2. Collective Memory:}
\begin{itemize}
\item Every educated Arab knows these dynasties
\item Taught extensively in schools across Arab world
\item Referenced in literature, poetry, media
\item Part of shared cultural heritage
\end{itemize}

\textbf{3. Symbolic Value:}
\begin{itemize}
\item Umayyad = Arab aristocracy, military expansion
\item Abbasid = Cultural golden age, intellectual flourishing
\item Both = Pre-modern governance models
\item Neither = Democratic or contemporary systems
\end{itemize}

\textbf{4. Rhetorical Power:}
\begin{itemize}
\item Instantly recognizable reference points
\item Carry weight of historical authority
\item Evoke specific imagery and associations
\item Enable sophisticated political critique
\end{itemize}

\textbf{The Nostalgia vs. Critique Tension:}

\begin{tabular}{|l|p{5cm}|p{5cm}|}
\hline
\textbf{Aspect} & \textbf{Nostalgic View} & \textbf{Critical View} \\
\hline
Umayyad/Abbasid & Golden age of Islamic civilization & Authoritarian dynastic rule \\
\hline
Modern reference & "Return to our glorious past" & "Escape from medieval thinking" \\
\hline
Governance & Strong unified leadership & Undemocratic hereditary power \\
\hline
Culture & Authentic Islamic values & Outdated social structures \\
\hline
Implication & We should emulate them & We must move beyond them \\
\hline
\end{tabular}

\textbf{The phrases in our analysis reflect the CRITICAL view, rejecting both nostalgia and the accusation of living in the past.}
\end{tcolorbox}

\clearpage

\section{Final Summary}

\begin{tcolorbox}[colback=boxcolor,colframe=accentcolor,breakable]
\textbf{Key Takeaways from This Phrase:}

\textbf{Grammatical Structures Mastered:}
\begin{itemize}
\item Negative copula \textarabic{لَيْسَ} and its complete conjugation
\item Use of \textarabic{لَيْسَ} with prepositional phrase predicates
\item Parallel structure with coordinating conjunctions
\item Case government: nominative subject + accusative predicate (with \textarabic{لَيْسَ})
\item Prepositional phrases functioning as predicates
\item Adjective agreement in gender, number, case, and definiteness
\item Difference between \textarabic{فِي} and \textarabic{بِ} prepositions
\item Full repetition vs. elliptical coordination
\end{itemize}

\textbf{Vocabulary Domains:}
\begin{itemize}
\item Negation patterns and particles
\item Historical dynasties and periods
\item Political and governmental terminology
\item Time and era vocabulary
\item Coordinating conjunctions
\item Nisba adjective usage
\end{itemize}

\textbf{Cultural Knowledge:}
\begin{itemize}
\item Umayyad and Abbasid historical context
\item Modern Arab political discourse patterns
\item Rhetorical use of historical references
\item Tension between tradition and modernity
\item Question-answer debate structures
\item Historical memory in contemporary politics
\end{itemize}

\textbf{Discourse Skills:}
\begin{itemize}
\item Understanding rhetorical questions vs. declarative statements
\item Recognizing parallel structure for emphasis
\item Analyzing tone and implication in political language
\item Comparing question and answer formulations
\item Identifying defensive vs. offensive rhetorical strategies
\end{itemize}

\textbf{Practical Application:}
\begin{itemize}
\item Expressing categorical negation in Arabic
\item Formulating emphatic denials
\item Engaging in political and historical discussions
\item Understanding contemporary Arab media discourse
\item Using historical references appropriately
\item Constructing parallel structures for rhetorical effect
\end{itemize}

\vspace{0.5em}
\textbf{Connection to Previous Phrase:}

This declarative statement serves as a direct response to the earlier interrogative:
\begin{itemize}
\item Question: \textarabic{هَلْ نَحْنُ بِعَصْرٍ عَبَّاسِيٍّ أَمْ عَصْرِ الدَّوْلَةِ الْأُمَوِيَّةِ؟}
\item Answer: \textarabic{لَسْنَا فِي الْعَصْرِ الْعَبَّاسِيِّ وَلَسْنَا فِي الدَّوْلَةِ الْأُمَوِيَّةِ}
\item Together they form a complete dialogue demonstrating question-answer patterns in Arabic political discourse
\end{itemize}
\end{tcolorbox}

\end{document}